\documentclass[tikz,border=10pt]{standalone}
\usepackage{tikz}
\usepackage{amsmath}
\usepackage{amssymb}
\usepackage{ctex}
\usetikzlibrary{calc, positioning, arrows.meta, decorations.pathreplacing, matrix}

% 对角化几何分解：x -> P^{-1}x -> Lambda P^{-1}x -> P Lambda P^{-1}x = Ax
% 用小网格变化示意四个步骤

\begin{document}
\begin{tikzpicture}[>=Stealth,
  stepbox/.style={draw=gray!60, fill=white, rounded corners=6pt,
                  minimum width=2.6cm, minimum height=2.6cm,
                  inner sep=5pt},
  arrowlabel/.style={font=\small, align=center},
  steplabel/.style={font=\small\bfseries, align=center},
]

%% ===== 步骤0：原始空间 x（标准基）=====
\begin{scope}[xshift=0cm, yshift=0cm]
  \node[stepbox] (box0) at (0,0) {};

  % 标准网格（单位正方形）
  \begin{scope}[xshift=-0.9cm, yshift=-0.9cm, scale=0.55]
    \draw[gray!30,thin] (0,0) grid (3,3);
    % 基向量 e1, e2
    \draw[->,blue!60!black,thick] (0,0) -- (1,0)
      node[below,font=\tiny] {$\mathbf{e}_1$};
    \draw[->,teal!60!black,thick] (0,0) -- (0,1)
      node[left,font=\tiny] {$\mathbf{e}_2$};
    % 向量 x
    \fill[black] (1.5,1.2) circle (2.5pt);
    \draw[->,orange,thick] (0,0) -- (1.5,1.2)
      node[right,font=\tiny] {$\mathbf{x}$};
  \end{scope}

  \node[steplabel, below=0.5cm of box0] {$\mathbf{x}$\\标准坐标};
\end{scope}

%% 箭头0→1
\node[arrowlabel, red!70!black] at (2.2, 0.5)
  {$P^{-1}$\\转到特征基};
\draw[->,very thick,red!70!black] (1.3, 0) -- (3.1, 0);

%% ===== 步骤1：特征基坐标 P^{-1}x =====
\begin{scope}[xshift=4.4cm, yshift=0cm]
  \node[stepbox] (box1) at (0,0) {};

  % 倾斜网格（特征向量方向）
  \begin{scope}[xshift=-0.9cm, yshift=-0.9cm, scale=0.55]
    % 特征向量 p1=(1,0), p2=(1,-1) 方向的斜网格
    \draw[gray!25,thin] (0,0) -- (3,0);
    \draw[gray!25,thin] (0,1) -- (3,1);
    \draw[gray!25,thin] (0,2) -- (3,2);
    \draw[gray!25,thin] (0,0) -- (1.5,3);
    \draw[gray!25,thin] (1,0) -- (2.5,3);
    \draw[gray!25,thin] (2,0) -- (3.5,3);
    % 特征基向量
    \draw[->,red!60!black,thick] (0,0) -- (1,0)
      node[below,font=\tiny] {$\mathbf{p}_1$};
    \draw[->,orange!70!black,thick] (0,0) -- (0.5,1)
      node[left,font=\tiny] {$\mathbf{p}_2$};
    % 向量 P^{-1}x
    \fill[black] (1.5,1.2) circle (2.5pt);
    \draw[->,orange,thick] (0,0) -- (1.5,1.2)
      node[right,font=\tiny] {$P^{-1}\mathbf{x}$};
  \end{scope}

  \node[steplabel, below=0.5cm of box1] {$P^{-1}\mathbf{x}$\\特征基坐标};
\end{scope}

%% 箭头1→2
\node[arrowlabel, green!50!black] at (6.6, 0.5)
  {$\Lambda$\\沿特征方向缩放};
\draw[->,very thick,green!50!black] (5.7, 0) -- (7.5, 0);

%% ===== 步骤2：缩放后 Lambda P^{-1}x =====
\begin{scope}[xshift=8.8cm, yshift=0cm]
  \node[stepbox] (box2) at (0,0) {};

  % 倾斜网格 + 拉伸
  \begin{scope}[xshift=-0.9cm, yshift=-0.9cm, scale=0.55]
    \draw[gray!25,thin] (0,0) -- (3,0);
    \draw[gray!25,thin] (0,1) -- (3,1);
    \draw[gray!25,thin] (0,2) -- (3,2);
    \draw[gray!25,thin] (0,0) -- (1.5,3);
    \draw[gray!25,thin] (1,0) -- (2.5,3);
    \draw[gray!25,thin] (2,0) -- (3.5,3);
    % 特征基向量（被拉伸：lambda1=3倍, lambda2=2倍）
    \draw[->,red!60!black,thick] (0,0) -- (3,0)
      node[below,font=\tiny] {$\lambda_1\mathbf{p}_1$};
    \draw[->,orange!70!black,thick] (0,0) -- (1.0,2)
      node[left,font=\tiny] {$\lambda_2\mathbf{p}_2$};
    % 缩放后向量
    \fill[black] (2.2,1.6) circle (2.5pt);
    \draw[->,orange,thick] (0,0) -- (2.2,1.6)
      node[right,font=\tiny] {$\Lambda P^{-1}\mathbf{x}$};
  \end{scope}

  % 对角矩阵提示
  \node[font=\tiny, gray!70, align=center] at (0.5, 0.8)
    {$\Lambda=\mathrm{diag}(\lambda_1,\lambda_2)$};

  \node[steplabel, below=0.5cm of box2] {$\Lambda P^{-1}\mathbf{x}$\\沿特征方向伸缩};
\end{scope}

%% 箭头2→3
\node[arrowlabel, blue!70!black] at (11.0, 0.5)
  {$P$\\转回标准基};
\draw[->,very thick,blue!70!black] (10.1, 0) -- (11.9, 0);

%% ===== 步骤3：结果 Ax（标准基）=====
\begin{scope}[xshift=13.2cm, yshift=0cm]
  \node[stepbox] (box3) at (0,0) {};

  % 标准网格（变换后）
  \begin{scope}[xshift=-0.9cm, yshift=-0.9cm, scale=0.55]
    \draw[gray!30,thin] (0,0) grid (3,3);
    \draw[->,blue!60!black,thick] (0,0) -- (1,0)
      node[below,font=\tiny] {$\mathbf{e}_1$};
    \draw[->,teal!60!black,thick] (0,0) -- (0,1)
      node[left,font=\tiny] {$\mathbf{e}_2$};
    % A*x (变换后的向量)
    \fill[red!70!black] (2.2,0.8) circle (2.5pt);
    \draw[->,red!70!black,thick] (0,0) -- (2.2,0.8)
      node[above,font=\tiny] {$A\mathbf{x}$};
  \end{scope}

  \node[steplabel, below=0.5cm of box3] {$A\mathbf{x}$\\回到标准坐标};
\end{scope}

%% ===== 底部公式框 =====
\node[font=\small, fill=white, draw=gray!50, thin, rounded corners=4pt,
      inner sep=10pt, align=center]
  at (6.6, -3.0) {
    对角化分解：$A = P\Lambda P^{-1}$\\[4pt]
    $A\mathbf{x}
      \;=\; \underbrace{P}_{\text{③转回}}
            \underbrace{\Lambda}_{\text{②缩放}}
            \underbrace{P^{-1}}_{\text{①转换}}\mathbf{x}$
  };

%% 整体标题
\node[font=\large\bfseries] at (6.6, 2.0)
  {对角化的几何分解：$A = P\Lambda P^{-1}$};

\end{tikzpicture}
\end{document}
